\section{La evaluación es el punto de llegada y también el de partida}

Al final de la entrevista se enfatiza que tanto la robótica como los grandes modelos enfrentan un cuello de botella en la evaluación. Sin una buena evaluación no es posible saber si el modelo de verdad se volvió más inteligente; sin una evaluación que pueda escalarse no es posible cerrar el ciclo de datos. La evaluación no es un informe que se redacta al terminar el entrenamiento, sino el punto de partida que determina la siguiente ronda de recolección de datos y la dirección del entrenamiento.

\subsection{Los tres papeles de la evaluación}

\noindent\begin{longtable}{p{0.22\textwidth}p{0.34\textwidth}p{0.35\textwidth}}
\toprule
Papel & Función & Ejemplos \\
\midrule
Árbitro & Juzgar si el modelo actual es más capaz & benchmark, tasa de éxito de extremo a extremo, tasa de aprobación en tareas reales. \\
Minero de datos & Exponer los fallos y los escenarios de cola larga & Los casos de fallo regresan de forma automática al conjunto de entrenamiento. \\
Volante de dirección & Decidir la siguiente dirección de entrenamiento y de recolección & Se agregan datos del tipo de tarea en el que la carencia es mayor. \\
\bottomrule
\end{longtable}

\begin{importantbox}{La evaluación también es un activo de datos}
Una evaluación de alta calidad define «qué cuenta como progreso». A medida que los modelos se vuelven más capaces, la evaluación misma se convierte en un dato escaso, porque su diseño requiere personas que entiendan mejor la tarea, los fallos y las necesidades futuras.
\end{importantbox}

\subsection{Resumen de la sección}

El ciclo de datos debe tener la evaluación como centro. Sin evaluación no hay dirección; sin dirección, los datos solo se acumulan y no se convierten en capacidad.
